\chapter{认知现象的结构解释}
\label{ch:interpretation}
\chapterguide{一个解释需要指出什么结构发生了变化，以及这种变化怎样影响可验证的任务表现。概念名称本身不承担证明。}

\section{无知、混淆与过度确信}
给定任务与实验族，无知可以表现为当前证据尚不足；结构性混淆则指两个目标冲突情境属于同一不可辨别类。前者可能通过继续当前实验改善，后者需要改变可取得或可使用的区别。

校准模型给出这一区别的可观察后果：只重复带偏差观察，会提高对$Y\oplus B$的确定性，但关于$Y$的Bayes错误仍为$1/2$。增加校准操作以后，目标可辨别，适当分配能够继续压低错误。因此“更加确信”需要绑定正在确信的变量，不能仅根据推理长度判断。

\section{好奇、专注与无效重复}
在覆盖模型中，选择实验可以针对尚未充分区分的冲突对。已经取得很多证据的方向继续占用大部分预算，会让其他方向成为最坏情境的瓶颈。切换操作具有明确的诊断理由。

这可以解释某些探索与专注行为：探索寻找未覆盖的错误方向，专注积累对当前关键方向的区别。它们的合理性依赖任务、已有证据与允许操作。覆盖率是固定全局分配的尺度，自适应阶段切换的最佳规则还需要更精细的分析。

无效重复可以用两种机制说明。第一种是结构性零增益：操作对当前关键冲突分布完全相同。第二种是资源性低增益：区别虽正，但相对于调用或读出成本很弱。将“无聊”解释成某种停止倾向时，应说明采用哪一种机制，避免仅把短期没有进展重新命名。

\section{理解与可迁移的组合}
在二元线性模型中，一份可验证的表示$l=a_1+\cdots+a_k$允许系统在全部未知情境上组合局部观察。它支持跨实例的同一操作方式。这样的迁移保证比记住某个实例的答案强，并可以通过不同情境下的风险检查。

但如果任务目标变为另一个$l'$，原表示未必仍然可用。理解的范围应由所覆盖的目标集合描述。例如，实验行空间给出全部可辨别线性目标；构造短表示的算法与它的成本，再给出实际可用的范围。

\section{发现、创新与新的可用区别}
一项改变至少可能发生在三个层次。发现新的实验，可能细分不可辨别类；在既有操作中找到更好的组织，可能改善诊断分配；找到低成本表示，可能保留信息同时减少取得答案的计算。三者各有不同的资源来源。

若要评价一个改变相对于旧系统的价值，可以固定任务族和预算，检验错误下限是否下降、达到指定可靠性所需调用是否减少、或计算与通信费用是否降低。新片段的表面差异不直接进入这些比较。

这里得到的是几种可分析的创新机制。一个具体变化是否属于其中之一，需要提供操作、概率或程序上的证据；“创新”这个名称无法代替这些结构。

\section{社会组织与成员差异}
两个成员使用不同措辞，但全部实验分布相同，它们的存在不会打破不可辨别类。成员使用互补实验，或一个成员能够把另一个成员难以解码的记录变成低成本充分消息，则可能改善整体系统。

对于同一完整信息与计算条件，一个集中控制者可以模拟有限团队。多成员优势需要说明额外资源：分散的观察权限、并行计算、局部专业读出，或者更易实现的模式。通信章的充分计数实现展示了一个明确情形：分散采样与集中似然判断可以得到相同的统计保证，同时只传递所需记录。

\begin{exercise}提出一个可反驳“更多独立成员必定消除共同错误”的有限实验，并说明独立性的条件是什么。\end{exercise}
\begin{exercise}对一个新的表示，分别设计信息保持、计算成本和跨实例可靠性三项检查。\end{exercise}
